% Slides 61--80
\sectionframe{04 / 3 April 2026}{The demo becomes a tested SOFR workstation}{Commits ca3d2d0, abfe66f, e13ff48, 7eaf671.}

\chronology{3 Apr 2026}{A missing SOFR fixing reveals the valuation boundary}
{Commit \texttt{ca3d2d0} fixes a runtime failure caused by a floating coupon requiring an unavailable historical fixing.}
{The live swap is moved to a two-business-day spot start.}
{A first Flask smoke-test file is added.}
{Schedule design and data availability must agree; avoiding accidental historical fixings is an explicit choice.}
{Commit ca3d2d0; tests/test\_smoke.py.}

\lesson{Spot starting avoids an accidental fixing dependency}
{A swap beginning on the evaluation date can contain an already-fixed or today's overnight observation.}
{The demo has no historical fixing store, so that coupon cannot be projected safely.}
{Advancing by two settlement business days makes all needed fixings prospective.}
{Production systems should ingest fixings; a sandbox may instead define a clean forward-start boundary.}
{portfolio.py at commit ca3d2d0.}

\chronology{3 Apr 2026}{SOFR curve construction is rebuilt around OIS}
{Commit \texttt{abfe66f} replaces generic annual deposits with an overnight front point and SOFR OIS helpers.}
{The curve becomes \texttt{PiecewiseLogCubicDiscount}.}
{Quoted calibration swaps are repriced and tested.}
{The market instrument must match the index and collateral economics represented by the curve.}
{Commit abfe66f.}

\lesson{The curve strip grows into a market object}
{Tenors ultimately span 1D, weekly, monthly, and annual pillars out to 30Y.}
{Each quote is normalized from percentage units into decimal rates.}
{The overnight point uses a deposit helper; longer points use \texttt{OISRateHelper} with \texttt{ql.Sofr}.}
{A single ordering contract links JSON quotes, labels, helpers, and diagnostics.}
{portfolio.py constants and build\_sofr\_curve.}

\begin{frame}{Bootstrap logic: solve one maturity at a time}
  \begin{enumerate}\setlength\itemsep{0.8em}
    \item Convert each market quote into the instrument helper that owns its conventions.
    \item Solve the next unknown curve segment so the new helper has zero NPV.
    \item Interpolate log discount factors between solved pillars.
    \item Reprice every input instrument and report its residual.
  \end{enumerate}
  \takeaway{The repricing table is the curve's unit test against its own market definition.}
  \src{portfolio.py build\_sofr\_curve and reprice\_sofr\_calibration\_swaps.}
\end{frame}

\lesson{Log-cubic discount interpolation encodes shape assumptions}
{Interpolation is performed on discount factors in logarithmic space, supporting positivity.}
{Cubic behavior produces a smoother curve than piecewise-linear zero rates.}
{Smoothness affects forwards and sensitivities between liquid pillars.}
{Interpolation choice is part of the model and should be recorded with the curve snapshot.}
{QuantLib PiecewiseLogCubicDiscount; portfolio.py.}

\lesson{Discount, zero, and forward views serve different users}
{Discount factors are consumed directly by pricing engines.}
{Continuous Actual/365 zero rates make the identity $P(0,T)=e^{-zT}$ transparent.}
{Daily one-day forwards reveal local interpolation behavior and scenario shape.}
{The workstation exposes all three rather than treating a curve as one line chart.}
{README.md; portfolio.py curve\_zero\_rate\_points and daily\_one\_day\_forward\_points.}

\lesson{Calibration residuals are expected near machine precision}
{Bootstrap helpers define the curve, so their model rates should match input quotes within numerical tolerance.}
{The repository reports market rate, fair rate, and NPV across pillars.}
{A residual spike often indicates a convention, ordering, or curve-range error.}
{Exact in-sample fit does not guarantee realistic interpolation or out-of-sample instruments.}
{build\_SOFR\_curve.py; tests/test\_sofr\_curve.py.}

\chronology{3 Apr 2026}{Repository standards and architecture are documented}
{Commit \texttt{e13ff48} adds architecture notes and Codex working rules.}
{The mission is defined as a small QuantLib sandbox with Flask and standalone scripts.}
{Quality gates become explicit commands rather than informal expectations.}
{Architecture documentation changes how later features are integrated and reviewed.}
{Commit e13ff48; AGENTS.md; docs/ARCHITECTURE.md.}

\lesson{The quality gates make the project reproducible}
{\texttt{unittest discover} checks the application and quantitative helpers.}
{\texttt{build\_SOFR\_curve.py} acts as an executable calibration report.}
{A Flask test-client request checks that the app can be created and routed.}
{A small project benefits from boring, copyable commands more than elaborate infrastructure.}
{AGENTS.md quality gates.}

\chronology{3 Apr 2026}{The rates workstation expands dramatically}
{Commit \texttt{7eaf671} adds roughly 2,547 lines and introduces the shared base template, dashboard, and trade form.}
{The market surface expands to a swaption normal-volatility matrix.}
{A Bermudan pricing grid and richer state management appear.}
{The repository shifts from a three-row demo to an interactive quantitative workstation.}
{Commit 7eaf671.}

\lesson{State normalization protects the pricing core}
{Defaults originate in JSON and Python constants, then user edits overlay selected fields.}
{Types, ranges, frequencies, directions, and labels are normalized before QuantLib construction.}
{The pricing functions receive one coherent state instead of raw form strings.}
{Normalization is the boundary between untrusted UI input and quantitative objects.}
{portfolio.py default\_portfolio\_state and normalize\_portfolio\_state.}

\lesson{The workstation carries a full ATM normal-vol matrix}
{Axes are defined explicitly by expiry and underlying-swap tenor labels.}
{Values are stored in basis-point normal volatility, then converted to decimal absolute volatility.}
{Interpolation supplies internal expiries needed by diagonal calibration without inventing on-screen pillars.}
{The UI shows quote provenance and the model shows which points each trade consumes.}
{README.md; portfolio.py swaption matrix helpers.}

\lesson{Normal volatility fits the rates regime}
{Bachelier pricing permits negative swap rates and quotes volatility in rate units.}
{A quote of 80 bp means an annualized absolute standard deviation of 0.0080.}
{QuantLib helpers must be told \texttt{ql.Normal}; passing the same number as lognormal vol changes meaning.}
{Unit conversion and volatility type are both contractual fields.}
{portfolio.py \_normal\_vol\_from\_bp and \_create\_swaption\_helper.}

\lesson{The dashboard becomes a diagnostic surface}
{It shows marks, market inputs, zero curves, forwards, OIS repricing, and a Bermudan grid.}
{Trade editors expose payment and reset frequencies rather than burying conventions.}
{Reset and scenario controls let users distinguish stored market state from perturbed state.}
{The UI is designed to explain pricing, not merely return a number.}
{templates/dashboard.html; templates/trade\_form.html.}

\lesson{A debug CSV creates an audit seam}
{Each calendar date can carry an input quote, zero rate, and one-day forward.}
{The file allows independent plotting and curve inspection outside Flask.}
{Paths are configured rather than exposed in the page response.}
{Small audit artifacts make numerical behavior easier to review and reproduce.}
{portfolio.py curve\_debug\_csv; routes.py curve debug download.}

\lesson{Curve risk becomes pointwise}
{Each quoted SOFR node is bumped and the trade repriced.}
{The resulting vector shows where value depends on the curve.}
{A swap should have zero vega cells; a swaption should respond to selected volatility pillars.}
{Risk-shape tests detect accidental market dependencies better than one aggregate NPV test.}
{portfolio.py trade\_point\_sensitivities; tests/test\_sofr\_curve.py.}

\lesson{A configurable valuation date replaces the machine clock}
{The default date becomes 10 March 2026 and anchors every schedule and curve.}
{User changes propagate through normalization and reset QuantLib's evaluation date.}
{Deterministic dates make screenshots, workbook comparisons, and tests reproducible.}
{A production request would carry an explicit as-of timestamp and market snapshot ID.}
{README.md; portfolio.py valuation\_date.}

\lesson{The curve phase establishes a reusable pattern}
{Build from appropriate helpers.}
{Reprice every input instrument.}
{Expose derived views and sensitivities.}
{The later USD and EUR XCCY curve bootstraps deliberately reuse this pattern.}
{portfolio.py and standalone\_xccy\_pricer.py build\_curves.}
